% Nouveau EGA V (cote 25): archive cote 25, site page 38 (PDF page 39).
% typescript. Draft compared: none.
% AI-transcribed by Claude Opus 5.5 (claude-opus-5-5, adaptive thinking, effort high), Message Batches API, batch ega25b (msgbatch_01A7n5YuXubUPNK1zsepi8RG), 2026-11-13; not yet checked by eye.
\scanpage{25/38}
\mspage{\hand{34}}

\var{\margin{$|$}}{\placedmargin{$|$}}{marginal note placed}{ega25_25_038-v0001}L'ensemble sous-jacent à ce préschéma n'est autre que l'ensemble des
points de $\underline{\overset{\hand{\unsure{/|\vee}}}{Y}}$= V($\hand{\phi}$) qui sont \uline{zéros singuliers}\var{\inserted[typed]{(par.16)}}{(par.16)}{insertion}{ega25_25_038-v0002} de la section $\hand{\phi}$ de
$\underline{O}_{X\mathrm{x}\overset{\hand{\vee}}{P}}$(1,1) \uline{relativement à la base} $\overset{\hand{\vee}}{P}$ , i.e. ses points \var{\xout{gé}}{}{typewriter strike-out}{ega25_25_038-v0003} à valeurs dans
un corps k sur $\overset{\hand{\vee}}{P}$ sont les points \var{\inserted[typed]{x}}{x}{insertion}{ega25_25_038-v0004} de $Y_k\hand{\subset} X_k$ tels que \var{\xout{\unsure{ou} bien $\hand{\phi}_k$ soit}}{}{typewriter strike-out}{ega25_25_038-v0005}
\var{\xout{\unsure{identiquement nulle en x, mais bien}}}{}{typewriter strike-out}{ega25_25_038-v0006} $\hand{\phi}_k$ s'annule à l'ordre deux au moins
en x, i.e. tel que \unsure{$W_k$|$M_k$} ne soit pas lisse de dimension relative r-1
sur k en x. La caractérisation énoncée des \var{\xout{points}}{}{typewriter strike-out}{ega25_25_038-v0007} zéros singuliers comm
les éléments du sousèpréschéma lisse P( $\hand{\underline{\nu}}_{X/P}$) de $X_{\overset{\hand{\vee}}{P}}$ provient aussitôt
de l'énoncé suivant, qui mérite de figurer en proposition préliminaire:
si S=Spec(k), et si H est un hyperplan de P , alors Y = X$\mathrm{x}_P$H est
\var{\xout{x \unsure{sur k} xxxxxxxxxxx}}{}{typewriter strike-out}{ega25_25_038-v0008}
\var{\xout{lisse en x le point rationnel sur k x si et seulement si}}{}{typewriter strike-out}{ega25_25_038-v0009} lisse sur k de
dimension relative d-1 en le point x $\hand{\in}$Y(k) (i.e. x est un zéro non
singulier, i.e. géom non singulier, de la section $\hand{\phi}$ de $\underline{O}_X$(1) définie
par H), si et seulement si H ne contient pas l'image par $\hand{\phi'}$ de l'espa-%
ce tangent à X en x (rel à k) - ou \var{\deleted{encore}}{}{deletion}{ega25_25_038-v0010} comme on dit\var{\xout{\illegible{1}}}{}{typewriter strike-out}{ega25_25_038-v0011} encore
(lorsque f:X$\hand{\longrightarrow}$P est une immersion, permettant d'identifier X à un
sous-préschéma de P) sss H \uline{n'est pas tangent à} X \uline{en} x. Cela résulte
trivialement du critère jacobien de lissité, ou de la définition d\unsure{u|a}
zéro singulier, une fois qu'on a bien précisé le sens de l'énoncé,
c'est à dire qu'on a précisé comment un \var{\xout{espace tangent}}{}{typewriter strike-out}{ega25_25_038-v0012} sous-espace
\var{\margin{$|$}}{\placedmargin{$|$}}{marginal note placed}{ega25_25_038-v0013}vectoriel de l'espace tangent à P en un point a (=f(x)) définit une
sous-variété linéaire de P (de telle façon que cela \var{\inserted[typed]{a}}{a}{insertion}{ega25_25_038-v0014} un sens de dire
que H ne contient pas ledit sous-espace vectoriel): cela provient
bien entendu de la suite exacte (b) ci-dessus, qui permet de\var{\xout{identifier}}{}{typewriter strike-out}{ega25_25_038-v0015}
définir une correspondance biunivoque entre l'ensemble des sous-espaces
vectoriels de l'espace tangent en a, et les sous-variétés linéaires de
P contenant a. Cette correspondance d'ailleurs revient à associer à

\var{\margin{\hand{38}}}{\placedmargin{\hand{38}}}{marginal note placed}{ega25_25_038-v0016}
